% ================================================================
% Chapter 13: Coalition Value and Totemic Solidarity
% ================================================================

\chapter{Coalition Value and Totemic Solidarity}
\label{ch:coalition}

\epigraph{The totality of beliefs and sentiments common to the average members
of a society forms a determinate system with a life of its own; one may call it
the \emph{collective} or \emph{common consciousness.}}{Émile Durkheim,
\textit{De la division du travail social} (1893)}

% ----------------------------------------------------------------
\section{Introduction}
% ----------------------------------------------------------------

Human groups are not mere collections of individuals---they are
\emph{coalitions} that compose their ideas jointly.  Durkheim's
\emph{mechanical solidarity} arises when members share a
\emph{conscience collective}: a jointly composed idea-system deeper than
any individual's contribution.  The totem, in Durkheim's analysis of
Australian religions, is precisely this: the symbol that expresses the
group's collective composition, the sign under which individual ideas are
subsumed into a shared cultural complex.\footnote{Durkheim,
\textit{Les Formes élémentaires de la vie religieuse} (1912).  The totem
is not merely a label but the material expression of the clan's
\emph{conscience collective}---what we formalize as the coalition composition.}

This chapter formalizes cooperative game theory in ideatic space.  We
introduce \emph{coalition composition}---the fold of a list of ideas via
the monoid operation, starting from void---and \emph{coalition value},
the depth of the joint composition.  We define \emph{marginal contribution},
the ideatic analogue of the Shapley marginal value, and prove eighteen
theorems showing that coalition formation is subadditive, void members
contribute nothing, associativity makes coalition merging well-defined, and
resonant coalitions produce resonant joint ideas.

The structural lesson is austere: there is no free lunch in cultural
synthesis.  The combined depth of a coalition never exceeds the sum of
individual depths.  The genius raises the group by at most their own
complexity; the free rider contributes nothing; and a coalition of voids
has zero cultural value regardless of its size.

% ----------------------------------------------------------------
\subsection{Historical Context: From Game Theory to Solidarity}
% ----------------------------------------------------------------

The mathematical foundations of coalition theory originate in John von~Neumann
and Oskar Morgenstern's \textit{Theory of Games and Economic Behavior}
(1944), which introduced cooperative game theory as a framework for analyzing
situations where binding agreements among players are possible.  Von~Neumann
and Morgenstern's key insight was that when agents can form coalitions, the
relevant strategic unit is no longer the individual but the \emph{group}:
the coalition's aggregate payoff determines each member's bargaining position.
The \emph{characteristic function} $v: 2^N \to \mathbb{R}$ assigns to each
subset~$S$ of the player set~$N$ a value $v(S)$ representing the total payoff
that coalition~$S$ can guarantee itself.  Our coalition value function is
the ideatic analogue: instead of monetary payoffs, it measures the cultural
depth that a group of idea-bearers can collectively achieve.\footnote{Von
Neumann and Morgenstern, \textit{Theory of Games and Economic Behavior}
(Princeton, 1944).  The cooperative game framework they developed has been
applied across economics, political science, and operations research; our
application to cultural composition is, to our knowledge, the first
systematic formalization in this direction.}

Lloyd Shapley's celebrated \emph{value} (1953) provides a unique allocation
of the coalition's total payoff to individual members, satisfying axioms of
efficiency, symmetry, linearity, and the null-player property.  The
\emph{Shapley value} of player~$i$ in a cooperative game with characteristic
function~$v$ is:
\[
  \phi_i(v) = \sum_{S \subseteq N \setminus \{i\}}
  \frac{|S|!\;(|N| - |S| - 1)!}{|N|!}
  \bigl[v(S \cup \{i\}) - v(S)\bigr].
\]
This formula averages the marginal contribution of player~$i$ across all
possible orderings in which the coalition might form.  Our marginal
contribution (Definition~\ref{def:marginal}) captures the same intuition in
the ideatic setting: the depth that an individual adds to a coalition depends
on when they join and what is already composed.  Shapley's null-player axiom
is our Theorem~\ref{thm:void-marginal}: a player who contributes nothing to
every coalition receives zero allocation.\footnote{Lloyd~S.\ Shapley, ``A
Value for $n$-Person Games,'' \textit{Annals of Mathematics Studies}~28
(1953), pp.\ 307--317.  For the connection between marginal contributions
and fair allocation, see also Alvin~E.\ Roth, ed., \textit{The Shapley
Value} (Cambridge, 1988).}

The \emph{core} of a cooperative game---the set of allocations under which
no coalition has an incentive to deviate---provides a complementary notion of
stability.  An allocation is in the core if and only if every coalition
receives at least its characteristic function value.  In ideatic terms,
a cultural arrangement is stable when every possible sub-group feels
adequately represented in the collective composition.  When the core is
empty, no stable arrangement exists, and the culture is prone to fission:
discontented sub-coalitions will break away to form their own collective
representations.  The subadditivity of our coalition value
(Theorem~\ref{thm:pair-subadd}) guarantees that the ``grand coalition''---the
entire society---always has a non-empty core, since no sub-coalition can
achieve more than the sum of its members' individual depths.

These game-theoretic tools find their anthropological counterpart in
Émile Durkheim's theory of solidarity.  In \textit{De la division du travail
social} (1893), Durkheim distinguished two ideal types of social cohesion.
\emph{Mechanical solidarity} characterizes segmentary societies---those
organized into structurally identical, self-sufficient units (clans,
lineages, moieties).  Here, solidarity arises from \emph{resemblance}:
members share a strong conscience collective, a common body of beliefs and
sentiments that integrates them through their very similarity.  The
collective representation is the fold of many similar contributions---our
coalition composition of ideas that largely resonate with one another.
\emph{Organic solidarity}, by contrast, characterizes complex societies
organized through the division of labor.  Here, solidarity arises from
\emph{complementarity}: members contribute different, specialized ideas
whose composition yields a collective representation richer than any
individual could produce alone.  The depth of the coalition exceeds any
singleton's depth precisely because the ideas are diverse.\footnote{Durkheim,
\textit{De la division du travail social} (Paris, 1893).  The distinction
between mechanical and organic solidarity maps onto our framework as follows:
mechanical solidarity corresponds to coalitions of resonant (similar) ideas,
while organic solidarity corresponds to coalitions of complementary
(non-resonant, high-depth) ideas.  Both are subadditive, but organic
solidarity tends to approach the subadditivity bound more closely.}

Durkheim's analysis of totemism in \textit{Les Formes élémentaires de la
vie religieuse} (1912) provides the ethnographic grounding for our
formalization.  Among the Australian Aboriginal peoples---particularly the
Aranda (Arrernte) of Central Australia and the Kariera of Western
Australia---social organization is structured by elaborate section systems
that partition the population into named groups, each associated with a
totemic emblem drawn from the natural world (kangaroo, emu, witchetty grub,
rain, etc.).  The Aranda system employs eight sections, arranged in a
pattern of prescribed marriage classes and moiety divisions.  The Kariera
system uses four sections.  In both cases, the totem functions as the
\emph{symbol} of the group's coalition composition: it is the material sign
under which individual contributions are subsumed into a collective
representation.  Durkheim argued that the totem is not merely a label or
taxonomic convenience but the very \emph{expression} of the group's social
force---the conscience collective made visible.\footnote{On Australian
section systems, see A.~P.~Elkin, \textit{The Australian Aborigines} (1938);
W.~Lloyd Warner, \textit{A Black Civilization} (1937); and more recently,
Patrick McConvell and Barry Alpher, ``The Formal Analysis of Kin Avoidance
Language in Australia,'' in \textit{Kinship Systems} (2002).}

The formal framework developed in this chapter also connects to contemporary
applications that Durkheim could not have anticipated.  Wikipedia, the
collaborative encyclopedia, exemplifies coalition composition at scale:
each editor contributes an idea (an edit, a paragraph, a citation) that is
composed with the existing article to produce a collective representation.
The article's ``depth''---its comprehensiveness, accuracy, and
sophistication---is the coalition value of the editor community.
Open-source software development follows a similar logic: each contributor's
code is composed with the existing codebase through version control, and the
project's complexity is the coalition value of the development team.  In
both cases, subadditivity holds: the project's depth never exceeds the sum
of all individual contributions, though it may fall far short of that
bound due to redundancy, conflicts, and the overhead of coordination.

Social capital theory, as developed by James Coleman, Robert Putnam, and
Pierre Bourdieu, provides a further bridge between cooperative game theory
and anthropological analysis.\footnote{James~S.\ Coleman, ``Social Capital
in the Creation of Human Capital,'' \textit{American Journal of Sociology}
94 (1988), pp.\ S95--S120; Robert~D.\ Putnam, \textit{Bowling Alone}
(2000); Pierre Bourdieu, ``The Forms of Capital,'' in J.~Richardson, ed.,
\textit{Handbook of Theory and Research for the Sociology of Education}
(1986).}  For Coleman, social capital inheres in the structure of
relations among persons---precisely the coalition structure we formalize.
For Putnam, social capital manifests as ``bridging'' (between diverse
groups) and ``bonding'' (within homogeneous groups), a distinction that maps
onto our framework: bonding capital corresponds to coalitions of resonant
ideas (mechanical solidarity), while bridging capital corresponds to
coalitions of diverse ideas (organic solidarity).  For Bourdieu, capital
exists in multiple forms---economic, cultural, social, symbolic---and is
convertible across forms through the habitus.  Our depth metric captures
one dimension of Bourdieu's cultural capital; the composition operation
captures the habitus's role in mediating between individual and collective
representations.

% ----------------------------------------------------------------
\section{Core Definitions}
\label{sec:ch13-defs}
% ----------------------------------------------------------------

\begin{definition}[Coalition Composition]
\label{def:coalition-compose}
The \emph{coalition composition} of a list of ideas $[a_1, a_2, \dots, a_n]$
is defined recursively:
\[
  \operatorname{coalitionCompose}([]) = \void, \qquad
  \operatorname{coalitionCompose}(a :: \text{rest}) =
    \operatorname{coalitionCompose}(\text{rest}) \compose a.
\]
This models a group's \emph{collective idea}---the joint composition of all
members' contributions.  In Durkheim's terms, this \emph{is} the conscience
collective.
\end{definition}

\begin{lstlisting}[caption={Coalition composition in Lean~4}]
def coalitionCompose : List I -> I
  | []         => void
  | a :: rest  => compose (coalitionCompose rest) a
\end{lstlisting}

Note the right-to-left fold: the earliest member in the list contributes
last, layered atop the previous members' composition.  This models the
anthropological fact that latecomers to a tradition interpret through---and
are interpreted against---the existing collective representation.

\begin{definition}[Coalition Value]
\label{def:coalition-value}
The \emph{coalition value} of a group is the depth of its joint composition:
\[
  \operatorname{coalitionValue}(\text{coalition}) =
  \depth(\operatorname{coalitionCompose}(\text{coalition})).
\]
\end{definition}

\begin{definition}[Marginal Contribution]
\label{def:marginal}
The \emph{marginal contribution} of idea $a$ to a coalition is the
difference in coalition value when $a$ is added:
\[
  \operatorname{marginal}(a, C) =
  \operatorname{coalitionValue}(a :: C) - \operatorname{coalitionValue}(C).
\]
This is the ideatic analogue of the Shapley marginal contribution in
cooperative game theory.
\end{definition}

% ----------------------------------------------------------------
\section{Empty and Singleton Coalitions}
\label{sec:ch13-empty}
% ----------------------------------------------------------------

\begin{theorem}[Empty Coalition is Void (13.1)]
\label{thm:empty-void}
The empty coalition composes to void: $\operatorname{coalitionCompose}([]) = \void$.
\end{theorem}

\begin{proof}[Proof sketch]
By definition of the base case.
\end{proof}

A group with no members has no culture.  This is the structural zero of
social organization---the ``state of nature'' in Hobbes's thought
experiment, but formalized as the void of ideatic composition.  Before the
first idea is contributed, the collective representation is empty.

\begin{theorem}[Empty Coalition Has Zero Value (13.2)]
\label{thm:empty-zero}
The cultural complexity of the empty group is zero:
$\operatorname{coalitionValue}([]) = 0$.
\end{theorem}

\begin{proof}[Proof sketch]
$\operatorname{coalitionCompose}([]) = \void$, and $\depth(\void) = 0$.
\end{proof}

\emph{Ex nihilo nihil fit}---from nothing, nothing comes.  A society begins
with zero cultural depth and must build it through composition.  The
Hobbesian ``war of all against all'' is not merely unpleasant; it is
culturally sterile.

\begin{theorem}[Singleton Coalition (13.3)]
\label{thm:singleton-coal}
A one-member coalition's composition is $\void \compose a$:
$\operatorname{coalitionCompose}([a]) = \void \compose a$.
\end{theorem}

\begin{proof}[Proof sketch]
By unfolding the recursive definition once.
\end{proof}

\begin{theorem}[Singleton Coalition Value (13.4)]
\label{thm:singleton-value}
A one-member coalition has value equal to that member's depth:
$\operatorname{coalitionValue}([a]) = \depth(a)$.
\end{theorem}

\begin{proof}[Proof sketch]
$\void \compose a = a$ by the void axiom, so
$\depth(\void \compose a) = \depth(a)$.
\end{proof}

The isolated individual contributes exactly their own depth to culture---no
more, no less.  Robinson Crusoe's cultural footprint is precisely his personal
complexity.  This theorem anchors the individual as the unit of cultural
measurement: before composition with others, each person's contribution is
exactly what they carry.\footnote{This should be contrasted with Durkheim's
anti-individualist position, which holds that the conscience collective is
\emph{emergent} and irreducible to individual contributions.  Our theorem
shows that at the singleton level, there is no emergence---emergence
requires at least two members.}

% ----------------------------------------------------------------
\section{Void Members and Inert Contributions}
\label{sec:ch13-void}
% ----------------------------------------------------------------

\begin{theorem}[Adding Void Preserves Coalition Composition (13.5)]
\label{thm:void-member}
Adding void to the front of a coalition does not change its composition:
\[
  \operatorname{coalitionCompose}(\void :: C) = \operatorname{coalitionCompose}(C).
\]
\end{theorem}

\begin{proof}[Proof sketch]
$\operatorname{coalitionCompose}(C) \compose \void =
\operatorname{coalitionCompose}(C)$ by the right-identity axiom.
\end{proof}

A ``member'' who brings nothing---the free rider, the passive observer---does
not alter the group's collective idea.  Durkheim's distinction between
mechanical and organic solidarity relies on each member \emph{contributing};
void members are structurally invisible.  This is the formal basis of
Olson's collective action problem: individuals who contribute nothing receive
the benefits of the collective good without increasing it.

\begin{theorem}[Void Member Has Zero Marginal Contribution (13.6)]
\label{thm:void-marginal}
Adding void to a coalition changes nothing about coalition value:
$\operatorname{marginal}(\void, C) = 0$.
\end{theorem}

\begin{proof}[Proof sketch]
Since $\operatorname{coalitionCompose}(\void :: C) =
\operatorname{coalitionCompose}(C)$ by Theorem~\ref{thm:void-member}, the
difference is zero.
\end{proof}

The free rider contributes nothing.  This formalizes Olson's central theorem
in \textit{The Logic of Collective Action}: non-contributors do not enhance
the collective good.\footnote{Mancur Olson, \textit{The Logic of Collective
Action} (1965).  Olson's argument that rational individuals will not
contribute to public goods unless compelled is captured by our result that
void contributions have zero marginal value.}

% ----------------------------------------------------------------
\section{Two-Member Coalitions}
\label{sec:ch13-dyad}
% ----------------------------------------------------------------

\begin{theorem}[Two-Member Coalition (13.7)]
\label{thm:two-member}
A two-member coalition $[a, b]$ has composition
$(\void \compose b) \compose a$.
\end{theorem}

\begin{proof}[Proof sketch]
By unfolding the definition twice.
\end{proof}

\begin{theorem}[Two-Member Coalition Simplified (13.8)]
\label{thm:two-member-simp}
The two-member coalition $[a, b]$ simplifies to $b \compose a$.
\end{theorem}

\begin{proof}[Proof sketch]
$\void \compose b = b$ by the void axiom, so
$(\void \compose b) \compose a = b \compose a$.
\end{proof}

The dyad---pair bond, marriage, master--apprentice---is the simplest
non-trivial coalition.  Its collective idea is simply the composition of
both members' ideas.  Simmel's insight that the dyad is qualitatively
different from larger groups is confirmed algebraically: the dyad is the
only coalition whose value reduces to a single composition, without nested
structure.\footnote{Georg Simmel, ``The Number of Members as Determining the
Sociological Form of the Group'' (1902).  Simmel argued that the dyad lacks
the super-individual structure that emerges in triads and larger groups.
Our simplification theorem confirms this: $[a, b]$ reduces to $b \compose a$,
a simple composition with no nesting.}

\begin{lstlisting}[caption={Two-member coalition simplified in Lean~4}]
theorem two_member_simplified (a b : I) :
    coalitionCompose [a, b] = compose b a := by
  simp [coalitionCompose, void_left]
\end{lstlisting}

% ----------------------------------------------------------------
\section{Subadditivity and Value Bounds}
\label{sec:ch13-bounds}
% ----------------------------------------------------------------

\begin{theorem}[Coalition Value Subadditivity for Pairs (13.9)]
\label{thm:pair-subadd}
The value of a two-member coalition is bounded by the sum of individual
depths:
\[
  \operatorname{coalitionValue}([a, b]) \;\le\; \depth(a) + \depth(b).
\]
\end{theorem}

\begin{proof}[Proof sketch]
By Theorem~\ref{thm:two-member-simp},
$\operatorname{coalitionValue}([a, b]) = \depth(b \compose a)
\le \depth(b) + \depth(a)$ by subadditivity.
\end{proof}

Two cultures merging cannot produce a joint culture deeper than the sum of
their parts.  Cultural synthesis is bounded---there is no ``$1 + 1 = 3$'' in
ideatic composition.  This constrains utopian visions of cultural fusion and
places a hard ceiling on the ``synergy'' of intercultural dialogue.  The
deepest possible synthesis of Shakespeare and Bashō is at most
$\depth(\text{Shakespeare}) + \depth(\text{Bashō})$.

\begin{example}[The Amish Barn Raising]
\label{ex:barn-raising}
The Amish barn raising exemplifies coalition composition in practice.
Each participant brings a specific skill---carpentry, roofing, masonry---and
the collective output is the composed structure.  The coalition value
(the completed barn) exceeds what any individual could achieve alone, yet
Theorem~\ref{thm:pair-subadd} ensures it cannot exceed the sum of all
individual depths.  The barn raising also illustrates
Theorem~\ref{thm:void-marginal}: community members who attend but
contribute nothing (void members) do not increase the coalition's cultural
output, though they may serve important social functions not captured by
the depth metric.

More precisely, let $a_1, a_2, \dots, a_n$ denote the ideas (skills,
techniques, knowledge) brought by the $n$ participants.  The barn as a
completed structure is $\operatorname{coalitionCompose}([a_1, \dots, a_n])$,
and its complexity---the coalition value---is the depth of this composition.
Subadditivity gives $\operatorname{coalitionValue}([a_1, \dots, a_n])
\le \sum_{i=1}^n \depth(a_i)$: the barn is at most as complex as the sum
of all individual skills.  In practice, the inequality is strict because
skills overlap (multiple carpenters), and the coordination overhead
(scheduling, communication) effectively reduces the realized depth below
the theoretical maximum.\footnote{Donald Kraybill, \textit{The Riddle of Amish
Culture} (2001), documents how barn raisings serve simultaneously as
economic cooperation and social bonding rituals.  The event reinforces
Durkheim's mechanical solidarity: participants share a common value system
(Ordnung) that makes their individual contributions resonant.}
\end{example}

\begin{remark}[Connection to Olson's Collective Action Problem]
\label{rem:olson}
Mancur Olson's \textit{The Logic of Collective Action} (1965) identifies
a fundamental tension in group behavior: rational self-interested agents
will under-contribute to public goods because they can free-ride on
others' contributions.  Theorem~\ref{thm:void-marginal} provides the
formal skeleton of this insight.  A free rider whose contribution is void
adds nothing to the coalition value, yet benefits from the collective
representation.  However, our framework also reveals a limit to Olson's
pessimism.  In ideatic space, the free rider is \emph{structurally
invisible}: they neither help nor harm the coalition's depth.  The
collective good---the coalition composition---is robust to the presence
of void members.  The problem is not that free riders destroy value but
that they fail to create it, leaving the coalition at a depth below its
potential maximum.  Olson's solution---selective incentives to encourage
participation---corresponds in our framework to mechanisms that convert
void members into non-void contributors, thereby increasing the realized
coalition value toward its theoretical bound.
\end{remark}

\begin{theorem}[Marginal Contribution Bounded by Depth (13.10)]
\label{thm:marginal-bounded}
Adding member $a$ to any coalition increases value by at most $\depth(a)$:
\[
  \operatorname{coalitionValue}(a :: C) \;\le\;
  \operatorname{coalitionValue}(C) + \depth(a).
\]
\end{theorem}

\begin{proof}[Proof sketch]
$\operatorname{coalitionCompose}(a :: C) =
\operatorname{coalitionCompose}(C) \compose a$, and
$\depth(\operatorname{coalitionCompose}(C) \compose a)
\le \depth(\operatorname{coalitionCompose}(C)) + \depth(a)$
by subadditivity.
\end{proof}

No individual can contribute more than their own depth to a group.  The
genius raises the group by at most their own complexity---the \emph{bounded
marginal returns} of individual talent in collective settings.  This is the
formal counterpart to the sociological observation that even extraordinary
individuals (Marx, Darwin, Freud) shift the collective culture by a
bounded amount proportional to their personal depth.

\begin{theorem}[Marginal Contribution at Most Depth (13.11)]
\label{thm:marginal-le-depth}
The marginal contribution of $a$ is at most $\depth(a)$:
\[
  \operatorname{marginal}(a, C) \;\le\; \depth(a).
\]
\end{theorem}

\begin{proof}[Proof sketch]
Since $\operatorname{coalitionValue}(a :: C) \le
\operatorname{coalitionValue}(C) + \depth(a)$ by
Theorem~\ref{thm:marginal-bounded}, the difference
$\operatorname{coalitionValue}(a :: C) - \operatorname{coalitionValue}(C)
\le \depth(a)$.
\end{proof}

An individual's marginal contribution to culture is bounded by their personal
depth.  No initiate brings more to a secret society than they carry.  This
constrains the ``great man'' theory of history: even the greatest individual's
marginal impact is bounded above.

\begin{example}[Wikipedia as Coalition Composition]
\label{ex:wikipedia}
Wikipedia's collaborative editing process provides a striking modern
illustration of coalition composition.  Each editor contributes an idea---a
paragraph of text, a citation, a structural reorganization, a factual
correction---that is composed with the existing article.  The article at
any given moment is $\operatorname{coalitionCompose}([e_1, e_2, \dots, e_n])$
where $e_i$ is the $i$-th edit in the article's history.  The article's
quality and comprehensiveness---its coalition value---is the depth of
this composition.

Several of our theorems find direct expression in Wikipedia's dynamics.
Theorem~\ref{thm:singleton-value} shows that a single editor's article
has exactly their own depth: one-person articles reflect one person's
knowledge.  Theorem~\ref{thm:pair-subadd} shows that merging two stub
articles produces an article no deeper than the sum of the stubs' depths.
Theorem~\ref{thm:void-marginal} explains why many minor edits (fixing
typos, adjusting formatting) do not substantially increase article depth:
these are near-void contributions whose marginal value is negligible.
And Theorem~\ref{thm:void-list} explains why articles created by bots
that contribute no substantive content remain at zero depth regardless
of their edit count.

The Wikipedia model also highlights a phenomenon not fully captured by
our current theorems: \emph{edit wars}, in which competing editors
repeatedly overwrite each other's contributions.  In our framework, an
edit war is a sequence of compositions and de-compositions (reversions)
that leaves the coalition value oscillating rather than monotonically
increasing.  The depth metric is indifferent to the social drama; it
tracks only the net compositional state.\footnote{For a sociological
analysis of Wikipedia's collaborative dynamics, see Yochai Benkler,
\textit{The Wealth of Networks} (2006), ch.~3; and Aaron Halfaker et al.,
``The Rise and Decline of an Open Collaboration System,''
\textit{American Behavioral Scientist} 57:5 (2013), pp.\ 664--688.}
\end{example}

\begin{remark}[Emergent Properties and the Limits of Subadditivity]
\label{rem:emergence}
Subadditivity (Theorem~\ref{thm:pair-subadd}) states that the coalition
value cannot \emph{exceed} the sum of individual depths.  But it says nothing
about how close the coalition value comes to that bound.  In practice,
coalition compositions exhibit a wide range of behaviors: some approach
the additive bound closely (when contributions are orthogonal and
non-redundant), while others fall far short (when contributions overlap or
interfere).

Durkheim's concept of \emph{emergence}---the idea that the collective
consciousness is qualitatively different from, and irreducible to, the sum
of individual consciousnesses---finds a precise location in our framework.
Emergence, in our terms, is the phenomenon whereby
$\operatorname{coalitionValue}(C) \ne \sum_{a \in C} \depth(a)$.  The
\emph{gap} between the coalition value and the additive bound is the
``emergence deficit''---the amount by which composition falls short of
mere aggregation.  When this deficit is small, the coalition is ``nearly
additive'' and emergence is weak.  When it is large, interactions between
members substantially reshape the collective idea, and emergence is strong.
The subadditivity theorem guarantees that emergence is always
\emph{non-positive}: coalitions never produce more than the sum of their
parts, though they may produce qualitatively different compositions that
are, in a structural sense, ``more than'' the mere concatenation of
individual ideas.
\end{remark}

% ----------------------------------------------------------------
\section{Associativity and Merging}
\label{sec:ch13-assoc}
% ----------------------------------------------------------------

\begin{theorem}[Coalition Composition Unfold (13.12)]
\label{thm:coal-unfold}
$\operatorname{coalitionCompose}(a :: \text{rest}) =
\operatorname{coalitionCompose}(\text{rest}) \compose a$.
\end{theorem}

\begin{proof}[Proof sketch]
This is definitional---the recursive case of \texttt{coalitionCompose}.
\end{proof}

Adding a new member to a coalition means composing their idea with the
existing collective idea.  Initiation rites formalize this structural
operation: the initiate's individual idea is composed with the group's
collective representation to produce a new, slightly altered collective.

\begin{theorem}[Three-Member Coalition (13.13)]
\label{thm:three-member}
The three-member coalition $[a, b, c]$ has composition
$((\void \compose c) \compose b) \compose a$.
\end{theorem}

\begin{proof}[Proof sketch]
By three applications of the recursive case.
\end{proof}

\begin{theorem}[Three-Member Simplified (13.14)]
\label{thm:three-member-simp}
The three-member coalition simplifies to $(c \compose b) \compose a$.
\end{theorem}

\begin{proof}[Proof sketch]
$\void \compose c = c$ by the void axiom, then unfold.
\end{proof}

The triad---Simmel's ``third''---introduces mediation.  The three-member
coalition's collective idea involves nested composition: the mediator
(middle member) connects the other two.  This is why Simmel considered the
triad qualitatively different from the dyad: the triad has internal
structure (nested composition) that the dyad lacks.  Formally, $[a, b]$
is a flat composition $b \compose a$, while $[a, b, c]$ is a nested
composition $(c \compose b) \compose a$---the nesting is the formal
signature of mediation.

\begin{example}[Jazz Improvisation as Coalition Composition]
\label{ex:jazz}
A jazz ensemble---say, the Miles Davis Quintet circa 1965, with Davis
(trumpet), Wayne Shorter (saxophone), Herbie Hancock (piano), Ron Carter
(bass), and Tony Williams (drums)---provides a vivid illustration of
coalition composition in real time.  Each musician contributes an idea
(a melodic line, a harmonic voicing, a rhythmic pattern) that is
composed with the ongoing collective sound.  The ensemble's music at
any moment is $\operatorname{coalitionCompose}([m_1, m_2, m_3, m_4, m_5])$
where $m_i$ is the $i$-th musician's contribution.

The triad theorem (Theorem~\ref{thm:three-member-simp}) is directly
relevant: any three musicians in the quintet form a trio whose collective
idea is a nested composition $(c \compose b) \compose a$.  The nesting
matters: the rhythm section (bass and drums) provides the compositional
foundation upon which the harmonic instrument (piano) layers, and the
soloists (trumpet, saxophone) compose atop this pre-existing structure.
The order of composition is not arbitrary---it reflects the functional
hierarchy of the ensemble.

Subadditivity (Theorem~\ref{thm:pair-subadd}) explains why even the
greatest ensemble cannot exceed the sum of its members' individual
virtuosity.  But the jazz ensemble also illustrates the \emph{emergence
deficit} discussed in Remark~\ref{rem:emergence}: the ensemble's music is
qualitatively different from a mere superposition of five solo
performances.  Interaction reshapes each contribution, producing a
coalition composition whose character cannot be predicted from individual
depths alone.  The depth of ``Kind of Blue'' is bounded by but not equal to
$\depth(\text{Davis}) + \depth(\text{Coltrane}) + \depth(\text{Evans}) +
\depth(\text{Chambers}) + \depth(\text{Cobb})$.\footnote{For a
musicological analysis of jazz ensemble interaction, see Ingrid Monson,
\textit{Saying Something: Jazz Improvisation and Interaction} (1996);
and Paul Berliner, \textit{Thinking in Jazz} (1994).}
\end{example}

% ----------------------------------------------------------------
\section{Resonance in Coalitions}
\label{sec:ch13-resonance}
% ----------------------------------------------------------------

\begin{theorem}[Self-Resonance of Coalition Value (13.15)]
\label{thm:coal-self-res}
Every coalition's collective idea resonates with itself:
\[
  \operatorname{coalitionCompose}(C) \res \operatorname{coalitionCompose}(C).
\]
\end{theorem}

\begin{proof}[Proof sketch]
By reflexivity of resonance.
\end{proof}

Every group's collective representation is self-consistent in the resonance
sense.  A culture never fails to ``recognize itself.''  Durkheim's
conscience collective is always internally coherent: the collective
representation may be shallow or deep, simple or complex, but it never
contradicts itself at the level of resonance.\footnote{This should not be
confused with logical consistency.  Resonance is weaker than logical
entailment: two ideas can resonate without being logically compatible.
The point is that the collective idea always \emph{recognizes} itself,
even if it contains contradictions by other metrics.}

\begin{theorem}[Resonant Members Yield Resonant Coalitions (13.16)]
\label{thm:resonant-singletons}
If $a \res a'$, then $\operatorname{coalitionCompose}([a]) \res
\operatorname{coalitionCompose}([a'])$.
\end{theorem}

\begin{proof}[Proof sketch]
Singleton coalitions simplify to the member itself (after void composition),
so the result follows from the hypothesis.
\end{proof}

If two individuals' ideas resonate---they ``get'' each other---their
one-person ``cultures'' are compatible.  This is the micro-foundation of
cultural affinity: interpersonal resonance at the individual level
translates directly to inter-cultural resonance at the singleton-coalition
level.

\begin{theorem}[Coalition Value of Void List (13.17)]
\label{thm:void-list}
A coalition consisting entirely of $n$ voids has zero value:
$\operatorname{coalitionValue}(\operatorname{replicate}(n, \void)) = 0$.
\end{theorem}

\begin{proof}[Proof sketch]
By induction on $n$: each void member preserves the existing composition
(Theorem~\ref{thm:void-member}), and the base case is the empty coalition
with value zero.
\end{proof}

A group of non-contributors produces no cultural depth---``a meeting of empty
minds.''  Bureaucracies without substance have zero cultural value regardless
of how many members they have.  The formal futility of scaling a null
organization: ten thousand voids compose to void.

\begin{lstlisting}[caption={Void coalition zero value in Lean~4}]
theorem void_coalition_zero_value : forall (n : ℕ),
    coalitionValue (List.replicate n void) = 0
  | 0     => by simp [coalitionValue, coalitionCompose, depth_void]
  | n + 1 => by
    simp [coalitionValue, coalitionCompose, List.replicate_succ, void_right]
    exact void_coalition_zero_value n
\end{lstlisting}

\begin{example}[Monastic Communities as Coalitions]
\label{ex:monastery}
The Benedictine monastery offers a counterpoint to the jazz ensemble:
a coalition designed to maximize collective depth through disciplined,
rule-governed composition.  The Rule of St.~Benedict prescribes a daily
schedule of prayer (\textit{opus Dei}), manual labor (\textit{labora}),
and sacred reading (\textit{lectio divina}) that structures each monk's
contribution to the communal life.  In our terms, each monk contributes
an idea shaped by the Rule, and the monastery's collective
representation---its spiritual culture---is the coalition composition
of these disciplined contributions.

Theorem~\ref{thm:void-list} is particularly illuminating here.  The
monastic tradition has always recognized the danger of \emph{acedia}---the
``noonday demon'' of spiritual torpor---which, in our framework,
corresponds to a monk whose contribution has become void.  A monastery
full of acedic monks has zero coalition value, regardless of its
population.  The elaborate monastic disciplines of self-examination,
confession, and spiritual direction are precisely mechanisms for detecting
and correcting void contributions before they accumulate.

The Benedictine emphasis on \textit{stabilitas loci}---the vow to remain
in one monastery for life---also finds formal expression.  By keeping the
coalition membership stable, the monastery ensures that each member's
contribution is composed with a familiar collective representation,
reducing the coordination overhead that plagues more fluid
coalitions.\footnote{On monastic organization as a form of social
engineering, see C.~H.~Lawrence, \textit{Medieval Monasticism} (3rd ed.,
2001); and Terrence Kardong, \textit{Benedict's Rule: A Translation and
Commentary} (1996).}
\end{example}

\begin{remark}[Connections to Chapter~\ref{ch:matching} and Matching Theory]
\label{rem:matching-connection}
The coalition theory developed in this chapter connects directly to the
matching theory of Chapter~\ref{ch:matching}.  There, we analyzed how
individual ideas are paired through the matching operation; here, we
study how those matched individuals compose into groups.  The connection
is precise: a coalition composition is a \emph{sequence} of pairwise
compositions, folded from right to left.  Every property of pairwise
composition (subadditivity, void identity, associativity) lifts to the
coalition level through induction on the list structure.

The matching-theoretic perspective suggests a natural refinement of
coalition value: instead of measuring only the depth of the grand
coalition's composition, one might also consider the depth profile of all
possible sub-coalitions---the \emph{characteristic function} of the
cooperative game.  Whether this richer structure yields additional
anthropological insights (e.g., about the stability of kinship systems
or the fission-fusion dynamics of band societies) remains an open question
for future work.
\end{remark}

\begin{theorem}[Coalition Value Monotonicity Bound (13.18)]
\label{thm:value-mono}
Adding any member to a coalition changes the value by at most the member's
depth:
\[
  \operatorname{coalitionValue}(a :: C) \;\le\;
  \operatorname{coalitionValue}(C) + \depth(a).
\]
\end{theorem}

\begin{proof}[Proof sketch]
This is Theorem~\ref{thm:marginal-bounded} restated.
\end{proof}

Cultural change is incremental.  Adding one person to a society can shift its
collective depth by at most that person's individual depth.  Revolutions
require deep individuals---shallow newcomers are structurally bounded in
their effect.  This is the formal counterpart to Kuhn's observation that
scientific revolutions require individuals of exceptional depth (Copernicus,
Newton, Einstein) to shift the collective paradigm.\footnote{Thomas Kuhn,
\textit{The Structure of Scientific Revolutions} (1962).  Kuhn's
``revolutionary scientist'' is, in our terms, an individual with
exceptionally high $\depth(a)$, whose marginal contribution is correspondingly
large---but still bounded.}

% ----------------------------------------------------------------
\section{The Anthropology of Collective Action}
\label{sec:ch13-collective-action}
% ----------------------------------------------------------------

The theorems of this chapter, taken together, provide a formal grammar for
the anthropological study of collective action.  We briefly survey how
this grammar illuminates several major themes in the anthropological
literature, drawing connections between the abstract algebra of coalition
composition and the concrete ethnography of human cooperation.

\subsection{Segmentary Lineage Systems}

Evans-Pritchard's classic analysis of the Nuer of southern Sudan (1940)
describes a \emph{segmentary lineage system} in which political coalitions
form and dissolve along genealogical lines.\footnote{E.~E.~Evans-Pritchard,
\textit{The Nuer} (Oxford, 1940).}  In the event of conflict, lineage
segments unite at successively higher levels of genealogical inclusion:
brothers unite against cousins, cousins against more distant kin, and
so on, following the principle ``I against my brother, my brother and I
against our cousin, our lineage against the world.''  In our framework,
each level of segmentary inclusion corresponds to a coalition composition
at a different scale.  The coalition value at each level is bounded by
subadditivity (Theorem~\ref{thm:pair-subadd}), and the marginal
contribution of each segment is bounded by its depth
(Theorem~\ref{thm:marginal-le-depth}).  The segmentary system is thus
a hierarchically nested sequence of coalition compositions, each
satisfying the structural constraints we have proven.

\subsection{Gift Exchange and Reciprocity}

Malinowski's analysis of the Kula ring (1922) and Mauss's comparative
theory of the gift (1925) both describe systems of reciprocal exchange
in which the act of giving creates social bonds---coalitions formed
through the exchange of culturally valued objects.\footnote{Bronisław
Malinowski, \textit{Argonauts of the Western Pacific} (1922); Marcel
Mauss, \textit{Essai sur le don} (1925).}  In Mauss's formulation, the
gift carries a portion of the giver's \emph{hau} (spiritual force),
which compels reciprocation.  In our terms, the gift is an idea with
positive depth, and giving it is an act of composition: the receiver's
cultural state is composed with the gift to produce a new state.  The
obligation to reciprocate ensures that the exchange is bidirectional,
producing a dyadic coalition (Theorem~\ref{thm:two-member-simp}) whose
value is at most the sum of both parties' depths.

\subsection{Collective Effervescence}

Durkheim's concept of \emph{collective effervescence}---the heightened
emotional energy generated by group ritual---maps onto our framework
as follows.  During collective ritual (the corroboree, the pilgrimage,
the political rally), individuals compose their ideas in rapid succession,
producing a coalition composition of exceptional density.  The emotional
intensity of the event reflects the rate of composition, not its ultimate
depth: subadditivity (Theorem~\ref{thm:pair-subadd}) still constrains
the final depth, but the subjective experience of rapid composition
creates the impression of transcendence.  Durkheim's observation that
participants feel themselves in the presence of a force ``greater than
themselves'' corresponds to the perception that the coalition value
exceeds any individual's depth---which it can, though it cannot exceed
the sum.

\subsection{Free Riding and Institutional Design}

The formal analysis of void members (Theorems~\ref{thm:void-member}
and~\ref{thm:void-marginal}) connects to a vast literature in
institutional economics on the design of mechanisms to prevent free
riding.  Elinor Ostrom's \textit{Governing the Commons} (1990) documents
how communities around the world have developed institutional rules---
monitoring, graduated sanctions, conflict-resolution mechanisms---to
ensure that members contribute to collective goods rather than free-riding
on others' contributions.\footnote{Elinor Ostrom, \textit{Governing the
Commons} (Cambridge, 1990).  Ostrom received the Nobel Prize in Economics
in 2009 for her work demonstrating that communities can self-organize to
manage common-pool resources without either privatization or state
regulation.}  In our framework, Ostrom's institutions are mechanisms for
converting void members into non-void contributors, thereby increasing
the coalition value.  The effectiveness of these institutions can be
measured by the gap between the \emph{actual} coalition value and the
\emph{counterfactual} value if all members contributed their full depth.

\subsection{Digital Communities and Online Collaboration}

The rise of digital platforms has created new forms of coalition
composition at unprecedented scale.  Open-source software projects
(Linux, Apache, Wikipedia) involve thousands of contributors whose
individual ideas are composed through version control systems and
collaborative editing interfaces.  The coalition value of these
projects---measured by code complexity, article comprehensiveness, or
feature richness---is the depth of the grand coalition's composition.
Our theorems apply directly: the project's depth is bounded by the sum
of all contributors' depths (subadditivity), non-contributing members
(lurkers, inactive accounts) add nothing (void marginal contribution),
and the project's depth grows at most linearly with the number of
substantive contributors (monotonicity bound).

These digital coalitions also exhibit phenomena that stretch our
framework in productive directions.  The ``1\% rule'' of online
communities---the empirical observation that roughly 1\% of users
create content, 9\% edit or curate, and 90\% passively consume---is a
natural consequence of Theorem~\ref{thm:void-list}: the 90\% of
passive consumers are void members whose presence does not increase the
coalition value.  Yet their presence may matter for other reasons
(audience size, network effects, advertising revenue) not captured by
the depth metric.\footnote{For the 1\% rule, see Ben McConnell and Jackie
Huba, ``The 1\% Rule,'' \textit{Church of the Customer Blog} (2006);
and Jakob Nielsen, ``Participation Inequality: Lurkers vs.\
Contributors in Internet Communities'' (2006).}

% ----------------------------------------------------------------
\section{Conclusion}
\label{sec:ch13-conclusion}
% ----------------------------------------------------------------

The cooperative game theory of ideatic space reveals a structural discipline
in coalition formation.  The empty coalition has zero value
(Theorem~\ref{thm:empty-zero}); singleton value equals individual depth
(Theorem~\ref{thm:singleton-value}); dyadic value is subadditive
(Theorem~\ref{thm:pair-subadd}); void members contribute nothing
(Theorem~\ref{thm:void-marginal}); and coalitions of voids are culturally
sterile (Theorem~\ref{thm:void-list}).

These results formalize Durkheim's insight that solidarity is not mere
aggregation but \emph{composition}---and composition is structurally
constrained.  The next chapter applies these coalition-theoretic tools to a
different institutional setting: the \emph{auction}, where ideas compete for
interpretive priority through depth-based bidding, connecting to the
anthropology of potlatch, prestige goods, and ritual economy.
